\cotefile{135}
\scanpage{3}
``Tordre'' une catégorie de Picard

1) Multiplier une Gr-cat par un groupe

Soit $C$ une (gr)-cat, $\Gamma$ un groupe, on va définir une nouvelle (Gr)-cat. notée $\Gamma\times C$. Le plus simple est de dire que $\Gamma$ désigne une catégorie discrète d'ens. sous-jacent $\Gamma$, et la loi de composition de $\Gamma$ définit un foncteur $\otimes$, avec une unique \deleted{structure}\inserted{contrainte} d'associativité. \deleted{\unsure{manière|Ceci|évidente}} [Il existe une \deleted{contrainte}\deleted{\unsure{contrainte}}\inserted{contrainte} de commutativité ssi $\Gamma$ est commutatif. (Les deux contraintes sont alors compatibles, et $\Gamma$ est une Catégorie \unsure{Abl|Abél|Ab|ACU})] On prend alors le produit des (Gr-cats) $\Gamma$ et $C$.
\[
\begin{aligned}
\pi_0(\Gamma\times C)&=\pi_0(\Gamma)\times\pi_0(C)=\Gamma\times\pi_0(C)\xrightarrow{\unsure{p|\pi}=\text{projection}}\pi_0(C)\\
\pi_1(\Gamma\times C)&=\pi_1(\Gamma)\times\pi_1(C)=\pi_1(C)\\
\unsure{k|K}_{\Gamma\times C}&=p^*(k_C)
\end{aligned}
\]

2) Supposons que $C$ soit $=$ une Pic-Cat, et \deleted{\unsure{une}}\inserted{$\Gamma$} commutatif. \deleted{Alors} $\Gamma\times C$ est une Pic-catégorie.

3) Si $C$ est une Pic-catégorie, les \inserted{autres} \deleted{structures}\inserted{contraintes} de comm. qui en font (avec sa \deleted{\unsure{contrainte}} d'associativité) une Pic-catégorie -- \unsure{i.e.}\deleted{\unsure{t.q.|types}} \deleted{\unsure{sont}} compatibles avec la contrainte d'ass., correspondent aux applications \uline{bilinéaires} \uline{antisymétriques}
\[
\varphi\colon\pi_0(C)\times\pi_0(C)\longrightarrow\pi_1(C)
\]
L'application correspondante
\[
\pi_0(C)\xrightarrow[\Sigma_\varphi]{}\pi_1(C)
\]
est donnée par
\[
\Sigma_\varphi(x)=\varphi(x,x)
\]
% archive sheet number: 2
